\begin{proposition}
\label{prop:ch15-nodal-local-envelope}
Let \(K\) be a complete nontrivially valued non-Archimedean field.  Choose
\(\varpi\in K\) with \(0<|\varpi|<1\), and put
\[
                  r:=|\varpi|,
          \qquad  a:=-\log r>0.
\]
Consider the reduced strict \(K\)-affinoid curve
\[
 U:=\mathcal M\bigl(K\langle x,y\rangle/(xy)\bigr),
 \qquad
 A:=V(y),\quad B:=V(x),\quad o:=V(x,y).
\]
Thus \(A\) and \(B\) are closed unit discs meeting only at \(o\).  Define
\(f:U\to\mathbf R\) by
\begin{equation}
\label{eq:ch15-nodal-obstacle}
 f(z):=
 \begin{cases}
  0,
       & z\in A,\\[2mm]
  -a\min\{r^{-1}|y(z)|,1\},
       & z\in B.
 \end{cases}
\end{equation}
Then \(f\) is continuous.  If
\[
 P_U^{\mathrm{raw}}f
 :=\sup\{g:g\text{ is locally Fubini--Study or semipositive
                    \(\mathbf R\)-PL on \(U\), and }g\leq f\},
\]
then
\begin{equation}
\label{eq:ch15-nodal-raw-envelope}
 P_U^{\mathrm{raw}}f(z)=
 \begin{cases}
  0,  &z\in A\setminus\{o\},\\
  -a, &z\in B,
 \end{cases}
\end{equation}
whereas
\begin{equation}
\label{eq:ch15-nodal-regularized-envelope}
 \bigl(P_U^{\mathrm{raw}}f\bigr)^*(z)=
 \begin{cases}
  0,  &z\in A,\\
  -a, &z\in B\setminus\{o\}.
 \end{cases}
\end{equation}
In particular, the upper-semicontinuous regularization is neither
classically psh nor CLD-v2 psh, and it is not continuous.
\end{proposition}

\begin{proof}
Since \(xy=0\), every point of \(U\) belongs to \(A\) or to \(B\), and
\(A\cap B=\{o\}\).  The two expressions in
\eqref{eq:ch15-nodal-obstacle} are continuous on their respective closed
branches and both take the value \(0\) at \(o\).  The finite closed-pasting
lemma therefore shows that \(f\) is continuous.

We first construct admissible minorants which detect the whole branch
\(A\setminus\{o\}\).  For \(m\geq1\), set, with the convention
\(\log0=-\infty\),
\begin{equation}
\label{eq:ch15-nodal-hm}
                  h_m:=\max\{m^{-1}\log|x|,-a\}.
\end{equation}
Although \(x\) vanishes identically on \(B\), the function \(h_m\) is a
real-valued \(\mathbf R\)-PL function.  Indeed, the two strict affinoid
domains
\[
 U_m^-:=\{|x|\leq r^m\},
 \qquad
 U_m^+:=\{|x|\geq r^m\}
\]
form a Laurent \(G\)-cover of \(U\), since \(r^m=|\varpi^m|\).  On \(U_m^-\)
one has
\[
                         h_m=-a=\log|\varpi|.
\]
On \(U_m^+\), the element \(x\) is an analytic unit: in the rational-domain
algebra its inverse is
\[
                  x^{-1}=\varpi^{-m}(\varpi^m/x).
\]
Consequently
\[
                         h_m=m^{-1}\log|x|
                         \qquad\text{on }U_m^+.
\]
These two \(\mathbf R\)-PL expressions agree on the overlap.  Notice in
particular that the zero locus of \(x\) is contained entirely in \(U_m^-\),
where the expression involving \(\log|x|\) has been replaced by the finite
constant \(-a\).

We next verify semipositivity, including at the node.  The normalization is
the finite surjective morphism
\[
 \nu:D_x\amalg D_y\longrightarrow U,
 \qquad
 D_x:=\mathcal M(K\langle x\rangle),
 \quad
 D_y:=\mathcal M(K\langle y\rangle).
\]
Its pullback is
\[
 h_m\circ\nu=
 \begin{cases}
  \max\{m^{-1}\log|x|,-a\},&\text{on }D_x,\\
  -a,&\text{on }D_y.
 \end{cases}
\]
On \(D_x\), this is subharmonic by the logarithm, positive-scaling, and
finite-maximum properties
\cite[Proposition~3.5(2),(4)]{GR25Classical}; on \(D_y\) it is constant and
hence subharmonic.  Subharmonicity descends under a finite surjective
morphism by \cite[Proposition~3.5(8)]{GR25Classical}.  Thus \(h_m\) is
subharmonic on the possibly singular curve \(U\).  On analytic curves,
subharmonicity agrees with classical plurisubharmonicity
\cite[\S4.1.2]{GR25Classical}.  Since \(h_m\) is real-valued
\(\mathbf R\)-PL, \cite[Theorem~4.3]{GR25Classical} now shows that \(h_m\)
is semipositive.

For \(z\in A\), one has \(|x(z)|\leq1\), so \(h_m(z)\leq0=f(z)\).
On \(B\), one has \(x=0\), hence \(h_m=-a\), while
\(-a\leq f\leq0\).  Therefore every \(h_m\) is admissible.  The constant
function \(-a=\log|\varpi|\) is also an admissible semipositive
\(\mathbf R\)-PL function.

It remains to show that no admissible function can exceed \(-a\) anywhere
on \(B\).  Put
\[
 D_r:=\{z\in B:|y(z)|\leq r\}
       \simeq\mathcal M\bigl(K\langle r^{-1}y\rangle\bigr).
\]
This closed disc has the unique Shilov point \(\eta_r\), namely the
\(r\)-Gauss point
\begin{equation}
\label{eq:ch15-nodal-r-gauss}
 \left|\sum_{n\geq0}c_ny^n\right|_{\eta_r}
       =\max_{n\geq0}|c_n|r^n.
\end{equation}
For completeness, after writing \(y=\varpi T\), the algebra of \(D_r\)
becomes \(K\langle T\rangle\).  The Gauss norm of every element of
\(K\langle T\rangle\) is its supremum norm on the closed unit disc and is
its value at the Gauss point.  Hence the singleton consisting of that point
is already a boundary for every analytic function.  Since the Shilov
boundary is the smallest nonempty closed boundary, it is precisely that
singleton.  Transporting it by \(y=\varpi T\) gives \(\eta_r\).

Every function occurring in the defining supremum is classically psh.  For
semipositive \(\mathbf R\)-PL functions this is
\cite[Theorem~4.3]{GR25Classical}.  Locally Fubini--Study functions are,
locally, positive scalar multiples of finite maxima of logarithms of
analytic functions, up to additive constants; hence their classical
plurisubharmonicity follows from
\cite[Proposition~4.2(1)--(3)]{GR25Classical}.  Classical psh functions are
functorial under analytic pullback, so an admissible \(g\) restricts to a
classically psh function on \(D_r\).  The affinoid maximum principle
\cite[Theorem~6.4]{GR25Classical} and the uniqueness of the Shilov point give
\begin{equation}
\label{eq:ch15-nodal-disc-maximum}
             g(z)\leq g(\eta_r)
                    \leq f(\eta_r)
                    =-a
             \qquad (z\in D_r).
\end{equation}
On \(B\setminus D_r=\{|y|>r\}\), the obstacle itself equals \(-a\).
Thus \(g\leq-a\) on all of \(B\).  Since the constant \(-a\) is admissible,
\begin{equation}
\label{eq:ch15-nodal-raw-on-b}
                         P_U^{\mathrm{raw}}f=-a
                         \quad\text{on }B.
\end{equation}

Now fix \(z\in A\setminus\{o\}\).  Then \(|x(z)|>0\), and hence
\[
 \lim_{m\to\infty}h_m(z)
 =\lim_{m\to\infty}
       \max\{m^{-1}\log|x(z)|,-a\}
 =0.
\]
Since every candidate is bounded above by \(f=0\) on \(A\), the raw
envelope equals \(0\) at \(z\).  Together with
\eqref{eq:ch15-nodal-raw-on-b}, this proves
\eqref{eq:ch15-nodal-raw-envelope}.

The sets
\[
 A\setminus\{o\}=\{|x|>0\},
 \qquad
 B\setminus\{o\}=\{|y|>0\}
\]
are open in \(U\).  Hence the raw envelope is locally constant away from
\(o\), and its upper-semicontinuous regularization has the same values
there.  Globally the raw envelope is at most \(0\).  On the other hand, the
rigid points
\[
                  z_n=(x=\varpi^n,y=0)\in A\setminus\{o\}
\]
converge to \(o\), and the raw envelope takes the value \(0\) at each
\(z_n\).  Therefore
\[
                 \bigl(P_U^{\mathrm{raw}}f\bigr)^*(o)=0,
\]
which proves \eqref{eq:ch15-nodal-regularized-envelope}.

Let \(u=(P_U^{\mathrm{raw}}f)^*\).  If \(u\) were classically psh, then its
restriction to \(D_r\) would be classically psh.  But
\[
                 \max_{D_r}u=u(o)=0,
 \qquad
                 u(\eta_r)=-a<0,
\]
contradicting the affinoid maximum principle
\cite[Theorem~6.4]{GR25Classical}.  Thus \(u\) is not classically psh.
Every CLD-v2 psh function is classically psh by
\Cref{thm:cld-to-classical-psh}; this is the implication recorded in
\cite[\S1.2]{GR25Classical}.  Hence \(u\) is not CLD-v2 psh either.
Finally, the points
\[
                  w_n=(x=0,y=\varpi^n)\in B\setminus\{o\}
\]
converge to \(o\), while \(u(w_n)=-a\) and \(u(o)=0\).  Therefore \(u\) is
not continuous.
\end{proof}
